
\section{The Federation of Free Congregations}

Editorial article in “Folkebladet” for January 26, 1898.

Those who have been so fortunate, we might almost say, as to have got hold of the Report of the Free Church’s Meeting in 1897, will, upon reading this valuable book, no doubt have come fully to see that such a federation of congregations as that of the Free Church is at any rate no freer—or, as the mockers say, "looser"—than the free connection among the apostolic congregations.

If we, who desire a free federation of congregations, are anarchists, then surely the Apostles must have been so as well. If it really is true that the New Testament is rule and guide, then the free federation of free congregations is at any rate not too free.

We have not yet seen anyone who, out of God’s Word, has been able to adduce a single statement that points in the direction that a Christian congregation is not absolutely independent and equipped with all the gifts and means that it needs for its existence, so that it can govern itself and has no need to be governed by any other power or authority.

That Christians love one another, and that Christian congregations love one another and support one another and help one another, that we know full well God’s Word teaches. But that they dominate one another and subdue one another, of that, for our part, we find not a trace.

Nothing therefore stands more firmly established, after the discussion that has been carried on since summer, than this: that a free federation of free congregations is biblical, apostolic, New Testament.

There is need of this also for the Free Church’s men, that under all the bitter mockery that is practiced against them and against their name and their work, they may have the firm and unshakable ground of God’s Word to stand upon.

We also intend to maintain—despite the learning of the learned—that the Lutheran Confession stands on the side of the Free Church in the question of the freedom and parity of the congregations. The Lutheran Church teaches that Lutheran congregations are joined together with one another into ecclesial unity, not by constitution, nor yet by ceremonies, but by the preaching or confession of the Gospel.

If anyone would claim that the free federation of congregations is anarchistic, then surely the Lutheran Confession must have been the same; for it also makes no claim that a majority dominion must be established in the Church.

Or perhaps there is someone who can show us that our Confession knows anything of annual meetings that rule over the congregations?

Georg Sverdrup
